\documentclass[CJK,aspectratio=169]{beamer}  %aspectratio控制页面的宽高比，16:9充满屏幕
%\documentclass[CJK,aspectratio=43]{beamer} %4:3宽高比

\usepackage{ctex} %中文字体设置
\usepackage{hyperref}
\usepackage{multicol}
\usepackage{makecell}

\usepackage{graphicx}
\usepackage{float}
\usepackage{subfigure}

\usepackage{amsmath}
\usepackage{mathtools}

\usepackage{soul}

\usepackage{diagbox}

\usepackage{tcolorbox}


% \usepackage[style=alphabetic,maxnames=3,minnames=1,maxbibnames=99]{biblatex}
% \addbibresource{ref.bib} 

\graphicspath{{figs/}}

\usetheme{Darmstadt}
% Darmstadt; Ilmenau
\usecolortheme{whale}
% \useoutertheme{sidebar}

% \usetheme[block=fill, sectionpage=none]{metropolis}
% \useoutertheme{infolines}
% \useinnertheme{metropolis}

\setlength{\parindent}{2em}

% \AtBeginSection[]
% {
% \begin{frame}
% \tableofcontents[currentsection]
% \end{frame}
% }

\begin{document}
	
	\title{Details on Parameterized Tests of Gravity Theories Using GWs with Matched Filtering}
	% \subtitle{子题目}
	\author[]{陈大文}
	\ctexset{today=old}
	
	\begin{frame}
		\vspace{0.5cm}
		\titlepage 
		\hypertarget{beginning}{}
	\end{frame}
	

	\begin{frame}[shrink,allowframebreaks] \transblindshorizontal<1>  %百叶窗效果
		\frametitle{\textsc{Contents}}
		\begin{multicols}{3}
			% \tableofcontents[hideallsubsections]
			\tableofcontents
		\end{multicols}
	\end{frame}


	% \AtBeginSection[]
	% {
	% \begin{frame}{\textsc{Contents}}
	% 	\tableofcontents[currentsection,currentsubsection,
	% 		hideothersubsections,
	% 		sectionstyle=show/shaded
	% 	]
	% \end{frame}
	% }

\section{检验类型}

\subsection{相容性检验}

\begin{frame}
	\frametitle{\textsc{Consistency Tests}}
	\par 不同gw检验间的关系：\href{https://arxiv.org/abs/2109.06988}{Investigating the relation between gravitational wave tests of general relativity} \cite{gw_tests_relation}

\end{frame}

\subsubsection{Residual Test}
\begin{frame}
	\frametitle{\textsc{Residual Test}}
	Residual Test的优点：
	\begin{itemize}
		\item[1. ] 不与双星参数简并\cite{test_GR_GW150914}
	\end{itemize}

\end{frame}

\subsubsection{IMR Test}


\subsubsection{QNM Test}


\subsubsection{Parameterized Tests}
\begin{frame}
	\frametitle{\textsc{Parameterized Tests}}
	
	\par Deviations 类型：
		\begin{table}[h!]
		\begin{center}
			\begin{tabular}{|c|c|}
				\hline
				\textsc{Type} & 表达式 \\
				\hline
				Fractional & $p_i \longrightarrow p_i(1 + \delta p_i)$ \\
				\hline
				Absolute/Additional & $p_i \longrightarrow p_i + \delta p_i$ \\
				\hline
				Particular Functional Form &  $p_i \longrightarrow p_i \cdot (1 + \delta p_i(\Xi;f))$\\
				\hline
			\end{tabular}
		\end{center}
	\end{table}

\end{frame}

\subsubsection{Dispersion Relation Test}


\subsubsection{Polarization Test}


\subsection{测量特定理论参数}

\begin{frame}
	\par 测量特定理论的参数，必须采用该理论全阶段non-GR波形，然而困境是\cite{wf_modelling_for_LISA}：
	\begin{itemize}
		\item [1. ] 理论结果；
		\item [2. ] 数值结果；
	\end{itemize}
	\par 折中方案是只研究特定stage，where 已有明确理论解。
	\par 另外的可能方案是\cite{ppE}。
\end{frame}

\section{GR波形模板误差}

\begin{frame}
	\frametitle{\textsc{Waveform Model Systematic Biases}}
	\begin{itemize}
		\item[1. ] GW150914\cite{test_GR_GW150914}: 
			\par SEOBNRv2\_ROM\_DOUBLESPIN: double-spin, non-precessing；
			\par IMRPhenomPv2: single-effective-spin, precessing；
			\par NR validation: we added the publicly available and new NR waveforms as mock signals to the data in the neighbourhood of GW150914 [36, 50, 51]；
			\par Detector Calibration: These uncertainties are modeled and included in the results presented here according to the treatment detailed in Ref. [3].
		\item[2. ]
	\end{itemize}
\end{frame}


\section{波形模板构造}

\subsection{参数化策略}
\subsubsection{Tiger}
\begin{frame}[allowframebreaks]
	\frametitle{\textsc{Tiger}}

\end{frame}

\subsubsection{ppE}
\begin{frame}[allowframebreaks]
	\frametitle{\textsc{ppE}}

\end{frame}

\subsubsection{FTI}
\begin{frame}[allowframebreaks]
	\frametitle{\textsc{FTI}}
	\par 针对SEOBNR模型的Inspiral阶段使用
\end{frame}

\subsubsection{pSEOBNR}
\begin{frame}[allowframebreaks]
	\frametitle{\textsc{pSEOBNR}}
	\par \cite{pPMR}
\end{frame}

\subsubsection{pyRing}
\begin{frame}[allowframebreaks]
	\frametitle{\textsc{pyRing}}
\end{frame}

\subsubsection{ParSpec}
\begin{frame}[allowframebreaks]
	\frametitle{\textsc{ParSpec}}
	
\end{frame}

\subsection{LALSimulation的IMRPhenom默认参数化}
\begin{frame}[allowframebreaks]
	\frametitle{\textsc{Parameterized PN Inspiral}}

	\begin{align}
		\Phi(f) = 2\pi ft_c - \varphi_c - \frac{\pi}{4} + \Psi(f_k)
	\end{align}
	\begin{align}
		\Psi(f) &= \text{CCE}(k;\eta,M;f) \cdot \sum_j \left[ \psi_j + \psi_{jl} \ln{f} \right] f^{j/3}\\
				&= \text{CCE}(k;\eta,M;f) \cdot \sum_j \left[ \alpha_j + \alpha_{jl} \ln{(2\pi Mf)} \right] (2\pi Mf)^{j/3}
	\end{align}
	\begin{align}
		\text{CCE}(\eta,M;f) = k \cdot \frac{3}{256}\frac{1}{\eta}(2\pi Mf)^{-5/3}
	\end{align}
	\par 其中 $f_k = \frac{f}{k}$；大多数情况下只考虑 $k = 2$ 的四极矩情况.

	\begin{itemize}
		\item[1. ] PN系数的定义选择：
		

		\begin{table}[h!]
			\begin{center}
				\begin{tabular}{|c|c|c|c|}
					\hline
					\diagbox{} & 幂项 & 公共系数 & PN系数 \\
					\hline
					\cite{reexpress_pn_3} & $(\pi Mf)^{j/3}$ & & $\psi_j,\psi_{jl}$ \\
					\hline
					IMRPhenomD & $(\pi Mf)^{j/3}$ & & $\alpha_j,\alpha_{jl}$\\
					\hline
					IMRPhenomXAS & $(\pi Mf)^{j/3}$ & & $\alpha_j,\alpha_{jl}$\\
					\hline
					LALSimulation & $(Mf)^{j/3}$ & & \makecell{幂项的$\pi^{1/3}$乘进系数，\\ 调整$\alpha_j,\alpha_{jl}$} \\
					\hline
					ppE\cite{ppE} & $(\pi\mathcal{M}_c f)^{j/3}$ & $\text{CCE}(\eta,\mathcal{M}_c;f)$ & \makecell{$M$变成$\mathcal{M}_c$，\\ 调整$\alpha_j,\alpha_{jl}$} \\
					\hline
					GW150914 \cite{test_GR_GW150914} & \multicolumn{3}{|c|}{IMRPhenomD} \\
					\hline
					GWTC-3 \cite{test_GR_GWTC3} & $\pi \tilde{f} = \pi \frac{GM}{c^3}(1+z)f$ & & \\
					\hline
				\end{tabular}
			\end{center}
		\end{table}
		\item[2. ] IMRPhenom波形family有显式的PN展开形式，
				而没有者、比如SEOBNR family，对应的是一种更 generic 的PN系数 test 方案：
		\begin{table}[h!]
			\begin{center}
				\begin{tabular}{|c|c|}
					\hline
					\textsc{Context} & PN系数的修改方法 \\
					\hline
					IMRPhenom & $\delta\varphi_j$  , Fractional or Additional(如$\delta\varphi_{-2}$，$\delta\varphi_1$) \\
					\hline
					SEOBNR & $\varphi_j \cdot \delta\varphi_j f^{(j-5)/3}$,log项同理\cite{test_GR_GW170817}\\
					\hline
				\end{tabular}
			\end{center}
		\end{table}
		\item[3. ] 是否将PN系数自旋依赖包括进偏离？
			\par 也就是考虑 $\varphi_j \longrightarrow (1 + \delta\varphi_j)\varphi_j^\text{NS} + \varphi_j^\text{S}$ \cite{test_GR_BBH_O1}；
			\par GWTC-3\cite{test_GR_GWTC3}中提到这种参数化可以避免系数偏离的 potentially singular behavior (奇异性原因来自non-spin和spin-dependent部分的 cancellation)。
		\item[4. ] 是否截断波形？
			\par 动力学情况不清楚的时候就要截断，比如用BNS事件GW170817 test GR\cite{test_GR_GW170817}。
	\end{itemize}
\end{frame}


\begin{frame}[allowframebreaks]
	\frametitle{\textsc{Parameterized Intermediate}}
	IMRPhenomXAS：
	\begin{align}
		\eta \varphi_\text{Int}^\prime = b_4f^{-4} + b_3f^{-3} + b_2f^{-2} + b_1f^{-1} + b_0 - \frac{4 c_0 a_\varphi}{(f-f_\text{RD})^2 + (2f_\text{damp})^2}
	\end{align}
\end{frame}

\begin{frame}[allowframebreaks]
	\frametitle{\textsc{Parameterized Ringdown}}
	IMRPhenomXAS (XAS中用 $p_i = -4, -2, -1/3$，D多一个$2$)：
	\begin{align}
		\eta \varphi_\text{RD}^\prime = c_\text{RD} + \sum_i c_i f^{-p_i} + \frac{c_0 a_\varphi}{f_\text{damp}^2 + (f-f_\text{RD})^2}
	\end{align}
\end{frame}



\subsection{Single/Multiple at a time}


\subsection{Hiearchical Model with Hyper-parameters}


\section{参数选择考量}

\subsection{Which Ones?}
\begin{frame}
	\frametitle{\textsc{Which Ones?}}
	\par 不正确的参数化造成的影响：\cite{false_parametrization}
\end{frame}

\subsection{SNR}
\subsubsection{Detector Sensitivity Band}
\subsubsection{事件IMR性质}

\subsection{物理现象映射}



\section{估计精确度}
\subsection{SNR}

\subsubsection{多事件}
\subsubsection{Detector Sensitivity Band}
\subsubsection{事件IMR性质}


\section{事件性质}
\begin{frame}[allowframebreaks]
	\frametitle{\textsc{Properties}}
	\begin{tcolorbox}
		\par \href{https://waveview.cardiffgravity.org/}{Cardiff: Gravitational Waveform Viewer}；
		\par \href{https://peviewer.igwn.org/}{IGWN: GW Event Viewer}；
		\par \href{https://gwosc.org/eventapi/html/GWTC/?page=1&pagesize=100}{GWOSC: GWTC}；
	\end{tcolorbox}

		\begin{table}[h!]
			\begin{center}
				\begin{tabular}{|c|c|}
					\hline
					事件 & 频带 \\
					\hline
					GW150914 & \makecell{$\dot{P}_\text{orb} \sim -0.1$ at $f_\text{GW} \sim 30 \text{Hz}$\\ $\dot{P}_\text{orb} \sim -1$ at $f_\text{GW} \sim 132 \text{Hz}$} \\
					\hline
					GW170817 & \\
					\hline
				\end{tabular}
			\end{center}
		\end{table}
\end{frame}


\section{不同模板的后验样本混合策略}


\section{GWs的约束结果}
\subsection{General Speakings}
\begin{frame}
	\frametitle{\textsc{General Speakings}}
	\begin{itemize}
		\item[$\bullet$] GW150914\cite{test_GR_GW150914}
			\par 强引力场动力学的约束不强，如一些PN系数的相对偏离 $\sim \mathcal{O}(1)$
			\par LIGO detectors are still a factor of a few away from their final design sensitivities.
	\end{itemize}
\end{frame}

\subsubsection{Residual Test}
\begin{frame}
	\frametitle{\textsc{Advantages}}
	% \begin{itemize}
	% \end{itemize}
\end{frame}

\subsubsection{IMR Test}
\begin{frame}
	\frametitle{\textsc{IMR Test}}
	\begin{itemize}
		\item[1. ] GW150914\cite{test_GR_GW150914}: 
			\par $f_\text{GW}^\text{end insp} = 132 \text{Hz}$;
			\par Expected $\text{SNR}_|text{det} \sim 19.5$, $\sim 16$ if truncated between $132$ and $1024$ Hz.
			\par We have verified that these conclusions are not a ected by the specifc formula [40, 60, 63] used to predict Mf and af, nor affected by the choice of fendinsp GW within 50Hz.
		\item[2. ]
	\end{itemize}
\end{frame}

\subsubsection{QNM Test}


\subsubsection{Parameterized Tests}


\subsubsection{Dispersion Relation Test}


\subsubsection{Polarization Test}


\subsection{测量特定理论参数}
\begin{frame}
	\frametitle{\textsc{IMR Test}}
	\begin{itemize}
		\item[1. ] GW150914\cite{test_GR_GW150914}: 
			\par $f_\text{GW}^\text{end insp} = 132 \text{Hz}$;
			\par Expected $\text{SNR}_|text{det} \sim 19.5$, $\sim 16$ if truncated between $132$ and $1024$ Hz.
			\par We have verified that these conclusions are not a ected by the specifc formula [40, 60, 63] used to predict Mf and af, nor affected by the choice of fendinsp GW within 50Hz.
		\item[2. ]
	\end{itemize}
\end{frame}


\section{\bf REFERENCES}
\begin{frame}[allowframebreaks]
	\frametitle{\textsc{REFERENCES}}
	% \cite{Mishra_2010}
	\bibliographystyle{plain}
	\bibliography{ref}
\end{frame}

\end{document}



